\originalchapter{来源与覆盖边界}
\label{ch:sources}

本卷的主来源为 Scott A. Hughes、MIT OpenCourseWare,
\href{https://ocw.mit.edu/courses/8-962-general-relativity-spring-2020/}{8.962 General Relativity, Spring 2020}.
理论以公开\href{https://ocw.mit.edu/courses/8-962-general-relativity-spring-2020/resources/mit8_962s20_guide/}{Lecture Guide}
和\href{https://ocw.mit.edu/courses/8-962-general-relativity-spring-2020/pages/lecture-summaries/}{Lecture Summaries}
第1--25讲的主题为范围, 重新组织为前九章;
并非课堂逐字翻译, 新增先修桥梁、推导与图均为本稿贡献.
\href{https://ocw.mit.edu/courses/8-962-general-relativity-spring-2020/pages/calendar/}{Calendar}
给出讲次和作业截止节点,
\href{https://ocw.mit.edu/courses/8-962-general-relativity-spring-2020/pages/syllabus/}{Syllabus}
给出原先修、课时、11套作业以及无考试的安排.

\originalsection{公开理论的边界}
第1--4讲建立洛伦兹不变量、四矢量、张量、体积元和守恒律;
第5--8讲讨论应力能量、等效原理、协变导数、平行输运、Lie导数及Killing对称;
第9--13讲连接测地线、曲率、潮汐、Bianchi恒等式和两种场方程推导;
第14--17讲进入线性引力、动态源和引力辐射;
第18--19讲讨论宇宙学、红移和距离;
第20--23讲讨论球对称恒星、TOV方程、紧致性、Schwarzschild黑洞与轨道.
Guide第24--25讲保留后牛顿展开、黑洞微扰、空间时间分解和数值相对论概要,
虽然这两讲在2020年未录视频, 仍是可引用的公开文字资料.
本卷为这些主题补充推导, 不能据此声称取得两讲全部课堂内容或掌握数值相对论全部方法.
Guide对Kerr及带电黑洞的简要介绍也不等于提供这些解的全部结构和证明.

各组讲次与本卷第1--9章的对应见卷首进度表.
第5讲按应力能量和联络分入前两章, 所以讲次在边界处有重叠.

第13讲的原Guide说明变分推导的教学更新因疫情未能按计划完成;
第15、17讲也提醒读者其材料较高级.
本卷对变分的边界项、规范和辐射能量等补足的推理属于独立重编,
没有把教学补充伪称为原教授逐字表述.
\href{https://ocw.mit.edu/courses/8-962-general-relativity-spring-2020/pages/readings/}{Readings}
中列出的教材只说明原课程背景, 本卷不捆绑其全文.

\originalsection{公开作业与缺题}
\href{https://ocw.mit.edu/courses/8-962-general-relativity-spring-2020/pages/assignments/}{Assignments}
明确说明: 2020年版作业不可获, 所列11套是Spring 2006版.
这与每份归档作业的第一页标题一致;
文件名中的S20和页尾OCW门户学期不能改变题目内文的年份.
第10--12章按原编号收入可获得的题面及本稿独立解答.

\begin{center}
\small
\begin{tabular}{@{}llll@{}}
\toprule
原作业 & 原编号数 & 本卷可恢复题面数 & 原题面缺口\\
\midrule
PS1 & 6 & 6 & 无\\
PS2 & 6 & 6 & 无\\
PS3 & 7 & 5 & 第5、6题\\
PS4 & 5 & 5 & 无\\
PS5 & 6 & 6 & 无\\
PS6 & 5 & 5 & 无\\
PS7 & 5 & 3 & 第4、5题\\
PS8 & 4 & 4 & 无\\
PS9 & 5 & 5 & 无\\
PS10 & 3 & 3 & 无\\
PS11 & 3 & 3 & 无\\
\midrule
合计 & 55 & 51 & 四项\\
\bottomrule
\end{tabular}
\end{center}

四个缺口分别为PS3.5的Carroll第3章第2题、PS3.6的第3章第3题,
PS7.4的第7章第3题和PS7.5的第7章第4题.
公开作业没有完整题面; PS3.5虽附旋度记号约定, 仍不足以恢复整个题目.
本卷保留编号、书目出处和缺口说明, 不凭记忆编题, 也不计入51道已有题面的大题.
2020阅读页列出的Carroll教材为2019出版版本, 而公开作业内文为2006;
旧题号与新版题号不能未经核对就等同.

PS7.3在作业PDF中只有Carroll第7章第1题的题号,
但\href{https://ocw.mit.edu/courses/8-962-general-relativity-spring-2020/resources/mit8_962s20_pset07_hint/}{官方PS7提示}
给出了完整拉氏量、构造线性化爱因斯坦张量的变分要求与对称性处理.
因此本卷恢复并解答的是该公开要求, 不是教材原页的逐字译文.
其余三份提示分别是
\href{https://ocw.mit.edu/courses/8-962-general-relativity-spring-2020/resources/mit8_962s20_pset01_hint/}{PS1第2题坐标含义澄清},
\href{https://ocw.mit.edu/courses/8-962-general-relativity-spring-2020/resources/mit8_962s20_pset03_hint/}{PS3第4题联络变换形式的等价说明},
\href{https://ocw.mit.edu/courses/8-962-general-relativity-spring-2020/resources/mit8_962s20_pset04_hint/}{PS4第2(a)题流体投影提示}.
这四份材料是局部指导, 不是完整官方解答集;
本卷所有作业解答均为独立推导, 对原题需补明坐标条件的地方在对应解答中校勘.
例如PS6.1对一般带shift坐标的约束分量作补明,
PS7.2区分原负交叉项的字面结果与改为正交叉项后满足原目标的相容版本;
原题的符号和要求没有被悄悄替换.
Syllabus明确无考试, 本卷不伪造期中或期末卷.

\originalsection{归档、许可与贡献}
选定原PDF为一份20页Guide、11套47页作业、四份8页提示, 共16份75页;
此外归档公开\href{https://ocw.mit.edu/courses/8-962-general-relativity-spring-2020/resources/ps10/}{PS10.nb}
作为原数值计算示例. Notebook仅静态审阅、未执行, 不计为额外答案,
阅读与重建本卷无需Mathematica.
源码保留原PDF和页尾OCW署名、网页快照、原网址与SHA256清单.

OCW选定资料和本卷中文改编按
\href{https://creativecommons.org/licenses/by-nc-sa/4.0/}{CC BY-NC-SA 4.0}
保留署名、非商业及相同许可要求.
MIT\href{https://ocw.mit.edu/pages/privacy-and-terms-of-use/}{使用条款}
与\href{https://mitocw.zendesk.com/hc/en-us/articles/4414756181403-How-is-all-rights-reserved-content-different-from-the-rest-of-OCW-content}{第三方材料FAQ}
是本轮许可判断依据.
对16份PDF的全部75页做了文字标记检查与渲染视觉查看,
未发现明确``All rights reserved''等保留权利标记.
PS5第2题的一幅矢量偏折示意未标单独图源或权利例外, 本卷采用自绘图.
没有发现标记不等于独立认证全部潜在权利; 逐对象证据保存在源码的\texttt{qa/source-audit.json}.

外链Hughes手写笔记因外站许可未确认而不打包.
商业教材全文、视频、照片和字体也不随本卷分发.
中文重编与题解为AI辅助独立贡献; MIT未审核或认可, 本卷只覆盖上述公开资料.
